% CS229 Final Project Poster -- 48 x 36 in landscape (beamerposter)
\documentclass[final]{beamer}
\usepackage[orientation=landscape,size=custom,width=121.92,height=91.44,scale=1.25]{beamerposter}
\usepackage[utf8]{inputenc}
\usepackage[T1]{fontenc}
\usepackage{lmodern}
\usepackage{amsmath,amssymb}
\usepackage{booktabs}
\usepackage{graphicx}
\usepackage{tikz}
\usetikzlibrary{arrows.meta,positioning}

% ----------------------------------------------------------------------
% Colors (Stanford Cardinal palette) + a lightweight block theme.
% ----------------------------------------------------------------------
\definecolor{cardinal}{RGB}{140,21,21}
\definecolor{coolgrey}{RGB}{77,79,83}
\definecolor{paloalto}{RGB}{0,86,63}
\definecolor{lightband}{RGB}{244,244,246}
\definecolor{good}{RGB}{0,110,70}
\definecolor{bad}{RGB}{170,30,30}

\setbeamercolor{block title}{fg=white,bg=cardinal}
\setbeamercolor{block body}{fg=black,bg=lightband}
\setbeamercolor{headline}{fg=white,bg=cardinal}
\setbeamerfont{block title}{size=\large,series=\bfseries}
\setbeamertemplate{navigation symbols}{}
\setbeamertemplate{caption}[numbered]
\setbeamercolor{caption name}{fg=cardinal}

% Tighter blocks with a little breathing room.
\setbeamertemplate{block begin}{
  \begin{beamercolorbox}[colsep*=0ex,dp=0.4ex,center]{block title}
    \vskip0.4ex\usebeamerfont{block title}\insertblocktitle\vskip0.3ex
  \end{beamercolorbox}
  \begin{beamercolorbox}[colsep*=0ex,vmode]{block body}
  \vskip0.6ex\begin{minipage}{0.96\textwidth}}
\setbeamertemplate{block end}{\end{minipage}\vskip0.6ex\end{beamercolorbox}\vskip1.2ex}

\newcommand{\figdir}{../figures}

% ----------------------------------------------------------------------
\begin{document}
\begin{frame}[t]

% ===================== HEADLINE BANNER =====================
\begin{beamercolorbox}[wd=\paperwidth,sep=0.6cm]{headline}
\centering
{\veryHuge\bfseries Micro-Transformer: A From-Scratch Autograd Engine}\\[0.3cm]
{\Huge\bfseries and a Gradient-Flow Ablation of Normalization \& Residual Connections}\\[0.5cm]
{\Large Manikanta Harsha Nadimpalli (\texttt{nmharsha}) \quad\textbullet\quad Luis Blanco Hoyos (\texttt{blancol})
 \quad\textbullet\quad CS229, General Machine Learning}
\end{beamercolorbox}
\vspace{0.8cm}

% ===================== THREE COLUMNS =====================
\begin{columns}[t]

% ---------------------------------------------------------------- COL 1
\begin{column}{0.315\paperwidth}

\begin{block}{The Question}
\large
Residual (skip) connections and LayerNorm are the scaffolding of every
transformer, usually treated as fixed. We make them the \textbf{object of study}:
\vskip0.4ex
\begin{center}\itshape
How do LayerNorm placement (pre / post / none) and residual connections
shape the gradient that reaches the \emph{early} layers during training?
\end{center}
\vskip0.4ex
\textbf{Input:} a window of character indices $x_{1:T}$.\quad
\textbf{Output:} next-char distribution $\hat p(\cdot\mid x_{1:t})$ at every $t$.\\
We watch not just the loss, but the intermediate gradient
$\partial L/\partial h_\ell$ at each block.
\end{block}

\begin{block}{Build It Ourselves $\rightarrow$ Measure Directly}
\large
A reverse-mode autograd engine in pure NumPy gives \emph{exact, low-overhead}
access to per-layer gradients --- no framework hooks to reverse-engineer.
\vskip0.6ex
\centering
\begin{tikzpicture}[
   node distance=0.55cm,
   box/.style={draw=cardinal,thick,rounded corners,fill=white,
               minimum width=10.5cm,minimum height=1.35cm,font=\large},
   tap/.style={draw=paloalto,thick,rounded corners,fill=paloalto!10,
               minimum width=10.5cm,minimum height=1.05cm,font=\large},
   ar/.style={-{Stealth[]},line width=1pt}]
  \node[box] (emb) {token emb $E$ \;$+$\; positional emb $P$};
  \node[box,below=of emb] (b1) {DecoderBlock $\times\,N$ \;(toggles: norm, residual)};
  \node[tap,below=of b1] (tap) {tap $h_\ell$ \;$\Rightarrow$\; record $\|\partial L/\partial h_\ell\|_2$};
  \node[box,below=of tap] (head) {final LayerNorm $\rightarrow$ linear head $\rightarrow$ logits};
  \draw[ar] (emb) -- (b1);
  \draw[ar] (b1) -- (tap);
  \draw[ar] (tap) -- (head);
\end{tikzpicture}
\vskip0.6ex
\raggedright\large
Validated to \textbf{machine precision} ($\sim\!10^{-10}$--$10^{-16}$) at every
layer vs.\ finite differences and a PyTorch oracle. \textbf{164 passing tests.}
PyTorch is used \emph{only} as a test oracle, never in the model.
\end{block}

\begin{block}{The Six-Cell Grid}
\large
3 norm placements $\times$ 2 residual settings, trained \emph{identically}
(same seed, data, AdamW, 2000 steps) so the only variable is the toggle.
\vskip0.4ex
\centering
\begin{tabular}{l|ccc}
\toprule
 & \textbf{pre} & \textbf{post} & \textbf{none} \\
\midrule
\textbf{residual on}  & $x{+}f(\mathrm{LN}\,x)$ & $\mathrm{LN}(x{+}f x)$ & $x{+}f x$ \\
\textbf{residual off} & $f(\mathrm{LN}\,x)$ & $\mathrm{LN}(f x)$ & $f x$ \\
\bottomrule
\end{tabular}
\end{block}

\end{column}

% ---------------------------------------------------------------- COL 2
\begin{column}{0.34\paperwidth}

\begin{block}{Hero Result: Per-Layer Gradient Norm}
\centering
\includegraphics[width=0.99\textwidth]{\figdir/gradflow_heatmap.pdf}
\vskip0.4ex
\raggedright\large
End-of-training gradient norm by depth ($\log_{10}$). \textbf{Residual-on}
(left 3 cols) stays $O(10^{-2})$ at every layer; \textbf{residual-off}
vanishes toward the input --- \texttt{none/no-res} is \emph{exactly zero} in
layers 0--2, so those layers never learn.
\end{block}

\begin{block}{Training Stability}
\centering
\includegraphics[width=0.92\textwidth]{\figdir/ablation_curves.pdf}
\vskip0.3ex
\raggedright\large
The three residual-on configs converge together near \textbf{1.8 nats};
the three residual-off configs plateau near \textbf{3.3}.
\end{block}

\begin{block}{Engine Validation Ladder}
\large
\begin{tabular}{ll}
\toprule
Component & Parity \\
\midrule
Every primitive op & finite diff, \texttt{atol 1e-7} \\
Attention backward & torch / closed form, $10^{-10}$ \\
LayerNorm & \texttt{torch.nn.LayerNorm}, $5{\times}10^{-16}$ \\
Decoder block & torch reconstruction, $10^{-10}$ \\
AdamW (25 steps) & \texttt{torch.optim.AdamW}, $10^{-10}$ \\
\bottomrule
\end{tabular}
\end{block}

\end{column}

% ---------------------------------------------------------------- COL 3
\begin{column}{0.315\paperwidth}

\begin{block}{Results: Residual Decides, Not Norm}
\large
\centering
\begin{tabular}{ll r r}
\toprule
Norm & Res & Val CE & Decay $\rho$ \\
\midrule
post & on  & \textbf{\textcolor{good}{1.775}} & 0.71 \\
pre  & on  & \textcolor{good}{1.844} & 0.68 \\
none & on  & \textcolor{good}{1.851} & 0.50 \\
\midrule
pre  & off & \textcolor{bad}{3.347} & 18.6 \\
post & off & \textcolor{bad}{3.347} & 30.7 \\
none & off & \textcolor{bad}{3.342} & 2746 \\
\bottomrule
\end{tabular}
\vskip0.5ex
\raggedright\large
Decay rate $\rho=(\|g_{\text{last}}\|/\|g_{\text{first}}\|)^{1/(L-1)}$:
$\rho\!\approx\!1$ preserves the gradient, $\rho\!\gg\!1$ vanishes it.
Baselines for scale: uniform 4.17, bigram floor 2.48, char-MLP 2.28 nats.
\end{block}

\begin{block}{The Takeaway}
\Large
\centering
At depth 4, the \textbf{\textcolor{cardinal}{residual connection is the
decisive factor}} --- not normalization placement.
\vskip0.5ex
\raggedright\large
\textbf{With} a skip, the gradient gets an identity path to the input and
stays $O(1)$ across depth; \textbf{all} norm placements then train and beat
the bigram floor. \textbf{Without} it, the gradient must cross every sublayer
Jacobian and decays geometrically --- early layers go untrainable regardless
of norm. Pre- vs.\ post-norm is a \emph{second-order} effect at this scale
(the known pre-norm advantage emerges deeper, without an adaptive optimizer).
\end{block}

\begin{block}{Scope, Limits \& Next Steps}
\large
Single seed/cell, single-head, one corpus, depth 4 --- so post-vs-pre is
``comparable,'' not a ranking. \textbf{Next:} multi-seed error bars, greater
depth (where pre/post diverges), multi-head, and Jacobian singular-value
spectra.
\end{block}

\begin{block}{Contributions \& Code}
\large
\textbf{Harsha:} data, baselines, scalar+tensor autograd, attention/LayerNorm/
block/transformer in-engine, AdamW, instrumentation, ablation, writeups.\\
\textbf{Luis:} closed-form attention-backward derivation + PyTorch oracle +
parity tests (the independent spec our engine is verified against).
\vskip0.4ex
\centering\large
\texttt{github.com/harshanm65/jacobigrad}
\end{block}

\end{column}

\end{columns}

% ===================== REFERENCES FOOTER =====================
\vspace{0.4cm}
\begin{beamercolorbox}[wd=\paperwidth,sep=0.35cm]{block body}
\footnotesize
\textbf{References:}\;
[1] Vaswani et al., \emph{Attention Is All You Need}, NeurIPS 2017.\quad
[2] He et al., \emph{Deep Residual Learning}, CVPR 2016.\quad
[3] Xiong et al., \emph{On Layer Normalization in the Transformer Architecture}, ICML 2020.\quad
[4] Ba et al., \emph{Layer Normalization}, arXiv:1607.06450, 2016.\quad
[5] Loshchilov \& Hutter, \emph{Decoupled Weight Decay Regularization}, ICLR 2019.\quad
[6] Karpathy, \emph{micrograd} \& \emph{char-rnn}, 2015/2020.
\end{beamercolorbox}
\end{frame}
\end{document}
